\chapter{The RC522 driver}
\label{ch:driver}

\chapterbrief
  {Turn ``shift bytes over SPI'' into ``read and write named registers''.}
  {Encapsulation, and composition (HAS-A).}
  {A real chip reporting \code{version() == 0x92}.}

\section{The problem}

Reading one RC522 register over raw SPI looks like this:

\begin{jcode}
byte[] tx = { (byte) (((0x37 << 1) & 0x7E) | 0x80), 0x00 };
byte[] rx = new byte[2];
link.transfer(tx, rx, 2);
int version = rx[1] & 0xFF;
\end{jcode}

Five lines, three magic numbers, and a bit-twiddle you must get exactly right.
Now picture that at every one of the roughly forty register accesses in a card
transaction: two hundred lines of identical incantation, each an opportunity to
type \hex{0x7F} where you meant \hex{0x7E}.

\section{The idea: gather the mess in one place}
\index{encapsulation}

That is what a class is, at its most useful: \textbf{a place to put the mess so
the rest of the program does not have to know about it.} After this chapter the
rest of your program says

\begin{jcode}
int version = driver.version();
driver.antennaOn();
\end{jcode}

and the address encoding exists in exactly one place, written once, correct once.

\begin{analogy}{The steering wheel}
A car's steering wheel is an interface to a mess: a rack, a pinion, hydraulic
assist, tie rods, geometry that changes with speed. You turn a wheel.

You are not permitted to reach into the steering column mid-corner and adjust
the rack --- not because it is secret, but because the mechanism only works if it
is driven through the one sanctioned control.

\code{private} is that housing. It is not about hiding secrets; it is about
\textbf{guaranteeing that the only way to affect this thing is through the paths
you have thought about.}
\end{analogy}

\subsection{The address encoding}

From the MFRC522 datasheet, the first SPI byte is an address byte:

\begin{center}
\begin{tikzpicture}[font=\scriptsize\sffamily, x=9.5mm, y=7mm]
  \foreach \i/\lbl in {0/7, 1/6, 2/5, 3/4, 4/3, 5/2, 6/1, 7/0}{
     \draw[draw=brewergrey] (\i,0) rectangle ++(1,1);
     \node[color=brewergrey, font=\tiny] at (\i+0.5, 1.28) {bit \lbl};
  }
  % RW bit
  \fill[brewerblue!16] (0,0) rectangle ++(1,1);
  \node at (0.5,0.5) {R/W};
  % address bits
  \fill[breweramber!16] (1,0) rectangle ++(6,1);
  \node at (4,0.5) {6-bit register address};
  % low bit
  \fill[black!7] (7,0) rectangle ++(1,1);
  \node at (7.5,0.5) {0};
  \draw[decorate, decoration={brace, amplitude=4pt, mirror}, brewergrey]
       (1,-0.12) -- (7,-0.12);
  \node[color=brewergrey, font=\tiny] at (4,-0.62) {\code{(addr << 1) \& 0x7E}};
  \draw[decorate, decoration={brace, amplitude=4pt, mirror}, brewergrey]
       (0,-0.12) -- (1,-0.12);
  \node[color=brewergrey, font=\tiny, align=center] at (0.5,-0.72) {1 = read\\0 = write};
\end{tikzpicture}
\end{center}

A register number must be shifted left one place to land in bits 6--1:

\begin{jcode}
((register << 1) & 0x7E) | 0x80    // read
 (register << 1) & 0x7E            // write
\end{jcode}

\code{\& 0x7E} is \code{0111 1110}: it masks off anything that overflowed past
bit 6 and forces bit 0 to zero. \code{| 0x80} sets the read flag. That is the
whole protocol, and after this chapter you never write it again.

\section{Build it}

\begin{jcode}
package com.midnightbrewer.hardware;

/**
 * The RC522 register-level driver.
 *
 * Encapsulation earns its keep here: the address encoding, the read/write bit
 * and the two-byte SPI dance all live IN HERE. The rest of the program calls
 * readRegister(0x37) and never writes ((reg << 1) & 0x7E) | 0x80 again.
 */
public class Rc522Driver {

    // ---- register addresses (MFRC522 datasheet, section 9) ----
    private static final int COMMAND_REG    = 0x01;  // CommandReg
    private static final int TX_CONTROL_REG = 0x14;  // TxControlReg (antenna)
    private static final int VERSION_REG    = 0x37;  // VersionReg (0x91 / 0x92)

    private static final int PCD_SOFT_RESET = 0x0F;  // the SoftReset command

    // ---- composition: the driver HAS-A SpiLink ----
    //
    // 'private' = nobody outside can reach past the driver and touch the bus.
    // 'final'   = once set in the constructor, this never points elsewhere.
    // The type is the INTERFACE, so the driver cannot tell real from fake.
    private final SpiLink link;

    // Dependency injection: the bus is handed in, not constructed here.
    public Rc522Driver(SpiLink link) {
        this.link = link;
    }

    // =================================================================
    // readRegister
    //
    //   C:  uchar Read_MFRC522(uchar addr) {
    //           SPI_Transfer( ((addr<<1)&0x7E) | 0x80 );  // address, read bit
    //           val = SPI_Transfer(0x00);                 // 2nd byte = answer
    //           return val;
    //       }
    //
    //   In Java ONE transfer() moves both bytes: build a 2-byte tx (encoded
    //   address, then 0x00), a 2-byte rx, call link.transfer, answer in rx[1].
    //   Return it masked with & 0xFF so 0x92 comes back as 146, not -110.
    // =================================================================
    public int readRegister(int register) throws SpiException {
        // TODO
        return 0;
    }

    // =================================================================
    // writeRegister
    //
    //   Same shape, two differences: the address byte does NOT set 0x80
    //   (that is what makes it a write), and the second byte is the value
    //   instead of 0x00. Nothing to read back.
    // =================================================================
    public void writeRegister(int register, int value) throws SpiException {
        // TODO
    }

    // =================================================================
    // setBitMask -- turn bits ON without disturbing the others.
    //
    //   C:  tmp = Read(reg);  Write(reg, tmp | mask);
    //
    //   "Read-modify-write". Build it from the two methods you just wrote,
    //   NOT from link.transfer. Each layer uses only the layer below it.
    // =================================================================
    public void setBitMask(int register, int mask) throws SpiException {
        // TODO
    }

    /** Turn bits OFF: same shape, but AND with ~mask. */
    public void clearBitMask(int register, int mask) throws SpiException {
        // TODO
    }

    /** Write the SoftReset command to CommandReg. One line. */
    public void softReset() throws SpiException {
        // TODO
    }

    // =================================================================
    // antennaOn -- power up the RF field.
    //
    //   C:  Read_MFRC522(TxControlReg);      // the C reads and ignores it
    //       SetBitMask(TxControlReg, 0x03);  // Tx1RFEn | Tx2RFEn
    //
    //   The two low bits of TxControlReg are the antenna drivers. Watch the
    //   register go 0x80 -> 0x83; that is read-modify-write doing its job.
    // =================================================================
    public void antennaOn() throws SpiException {
        // TODO
    }

    /** The chip's version byte: 0x91 (v1.0) or 0x92 (v2.0). */
    public int version() throws SpiException {
        // TODO
        return 0;
    }
}
\end{jcode}

\begin{concept}{Each layer uses only the layer beneath it}
\code{setBitMask} could call \code{link.transfer} directly and it would work.
Do not.

If you shortcut a layer you have created a second path to the same hardware ---
and the day you need to add logging, retries, or a mutex to register access, one
path gets it and the other silently does not. That is not a hypothetical; it is
precisely how ``I fixed the bug but it still happens sometimes'' bugs are born.
\end{concept}

\section{Verify --- with no hardware}

\begin{jcode}
package com.midnightbrewer.hardware;

import org.junit.jupiter.api.Test;
import static org.junit.jupiter.api.Assertions.assertEquals;

class Rc522DriverTest {

    @Test
    void readsTheVersionRegister() throws SpiException {
        FakeSpiLink link = new FakeSpiLink();
        link.setRegister(0x37, 0x92);

        Rc522Driver driver = new Rc522Driver(link);

        assertEquals(0x92, driver.version());
    }

    @Test
    void writeEncodesTheAddressWithoutTheReadBit() throws SpiException {
        FakeSpiLink link = new FakeSpiLink();
        Rc522Driver driver = new Rc522Driver(link);

        driver.writeRegister(0x14, 0x03);

        byte[] tx = link.sent.get(0);
        assertEquals((0x14 << 1) & 0x7E, tx[0] & 0xFF);   // read bit CLEAR
        assertEquals(0x03, tx[1] & 0xFF);
    }
}
\end{jcode}

The second test is the interesting one: it asserts on \textbf{what went out on
the wire}, which is the only way to catch an encoding bug. A test that checked
only return values would happily pass with the read bit set on a write.

\begin{verifybox}
\begin{shell}
mvn test -Dtest=Rc522DriverTest
\end{shell}
\end{verifybox}

\section{Now the real chip}

\begin{jcode}
package com.midnightbrewer.hardware;

public final class HardwareCheck {

    private HardwareCheck() { }

    public static void main(String[] args) throws Exception {
        System.out.println("Opening the reader...");

        try (Pi4jSpiLink link = new Pi4jSpiLink()) {
            Rc522Driver driver = new Rc522Driver(link);

            driver.softReset();
            Thread.sleep(50);              // the chip needs a moment after reset
            driver.antennaOn();

            int version = driver.version();
            System.out.printf("version() = 0x%02X%n", version);
            if (version == 0x91 || version == 0x92) {
                System.out.println("  -> the driver works on silicon.");
            } else {
                System.out.println("  -> chip not answering. Check wiring/power.");
            }
        }
    }
}
\end{jcode}

\begin{shell}
# on your laptop
rsync -av --exclude target/ ./ martin@raspberrypi.local:~/midnight-brewer/

# on the Pi
cd ~/midnight-brewer
mvn -q -DskipTests compile dependency:copy-dependencies -DincludeScope=runtime
java -cp "target/classes:target/dependency/*" \
     com.midnightbrewer.hardware.HardwareCheck
\end{shell}

\begin{verifybox}
\begin{trace}
version() = 0x92
  -> the driver works on silicon.
\end{trace}
\end{verifybox}

\subsection{If you get \texttt{0x00} or \texttt{0xFF}}

Those are the two ``nobody home'' values, and they are diagnostic:

\begin{itemize}
  \item \hex{0x00} --- MISO stuck low. Usually MISO not connected, chip
        unpowered, or CS on the wrong pin.
  \item \hex{0xFF} --- MISO stuck high, i.e.\ floating. Usually MISO connected
        to nothing, or a dead chip.
\end{itemize}

Check power first (3.3\,V, \emph{not} 5\,V), then wiring, then shorten your
wires. Do not debug the code: it is ten lines and your test already proved it.

\begin{deeper}{try-with-resources, and why \texttt{AutoCloseable} mattered}
\code{try (Pi4jSpiLink link = ...)} calls \code{close()} when the block exits ---
normally, by exception, or by early return. That is why \code{SpiLink extends
AutoCloseable} back in Chapter~\ref{ch:spi}.

On a kiosk it means a crash can never leave the RF field energised, which is
the sort of thing that quietly cooks a reader module overnight. Resource cleanup
that depends on remembering to write it is resource cleanup that eventually does
not happen.
\end{deeper}

\section{Traps}

\begin{traptable}
\code{version()} $= -110$ & Missing \code{\& 0xFF}. \\
\code{version()} $=$ \hex{0x00} & MISO low: power, wiring, or wrong CS pin. \\
\code{version()} $=$ \hex{0xFF} & MISO floating: not connected. \\
Reads work, writes do not & Read bit \hex{0x80} accidentally set on writes. \\
Fine on the bench, flaky when mounted & Wire length. Shorten them. \\
\hex{0x92} in the test, \hex{0x00} on the Pi & Stale \code{target/}. Run \code{mvn clean compile}. \\
\end{traptable}

\section{Commit}

\begin{shell}
git add -A
git commit -m "M2: Rc522Driver -- register read/write, bitmasks, reset, antenna"
\end{shell}
